\appendix

\section{Theoretical Motivation and Empirical Validation}
\label{appendix:spectral_separation}

To provide theoretical context for our router's design, we discuss the relationship between the semantic structure in Figure~\ref{fig:lmsys_holdout_structure} and the expected convergence behavior of the meta-algorithm.

In contextual bandit theory, the probability of a routing error (selecting a sub-optimal model) is inversely proportional to the margin of separation in the feature space. Even with only $3.10\%$ variance in PC1, the cluster structure creates a natural separation that aids the decision engine.

\subsection{The Margin of Decision}

We define the decision boundary $H$ at the unsupervised silhouette-optimal threshold $\text{PC1} = 0.222$. For any prompt $x$, the distance to this boundary is $d(x, H)$. Routing is facilitated because the density gap between clusters ensures that:
\begin{equation}
P(d(x, H) < \epsilon) \approx 0
\end{equation}
for a sufficiently small $\epsilon$. This means that relatively few prompts sit in the ``ambiguous zone,'' allowing the router to achieve stable model selection.

\subsection{Expected Routing Regret}

Because the clusters are spatially distinct ($n=606$ Low PC1 vs.\ $n=144$ High PC1 on the held-out set), the Thompson Sampling mechanism is expected to converge at a rate determined by the sub-Gaussianity of the rewards within each cluster. For standard Thompson Sampling with a spectral gap $\Delta_{\text{gap}}$ between the Low and High PC1 clusters, the theoretical regret bound takes the form:
\begin{equation}
R(T) \le \sum_{a \in \mathcal{A}} \frac{8 \ln(T)}{\Delta_{\text{gap}}} + \left(1 + \frac{\pi^2}{3}\right) \sum_{a \in \mathcal{A}} \Delta_a
\end{equation}
where $\mathcal{A}$ is the set of available models and $\Delta_a$ is the sub-optimality gap for model $a$.

Since PC1 shows significant rank correlation with model preference ($\rho = -0.370$, Figure~\ref{fig:lmsys_holdout_structure}), the feature space provides meaningful signal for the bandit to exploit. While our implementation includes practical modifications (fixed learning rates, selective expert updates), the empirical results in Section~\ref{sec:results} demonstrate that the system achieves near-optimal performance consistent with this theoretical intuition.
